% !TEX root = ../main.tex
%% Réussir un devoir surveillé : méthode, rédaction, capacités expérimentales et fiches expérimentales

\editionPage{methodes-devoir}{Réussir un devoir surveillé}{Méthode et organisation}{Deux heures, trois exercices · Comprendre ce qui est demandé}
\begin{blocEdition}{methodes-devoir}
\fontsize{10}{13.5}\selectfont
\editionEtape{Le devoir en bref}
\begin{tabularx}{\linewidth}{@{}p{38mm}Y@{}}\toprule
Durée et barème & 2 heures ; trois exercices indépendants, notés au total sur 20 points ; le barème de chaque exercice est
indiqué en tête.\\[1mm]
Programme & Plusieurs chapitres consécutifs ; les notions des chapitres précédents peuvent servir d'outils.\\[1mm]
Documents & Textes, schémas, graphiques, tableaux de mesures, programmes en Python : tout ce qui est utile est fourni,
mais il faut aller le chercher.\\[1mm]
Calculatrice & Autorisée ; elle sert aussi aux statistiques (moyenne, écart-type).\\\bottomrule
\end{tabularx}
\editionEtape{Gérer les deux heures}
\begin{questions}
\item Lire tout le sujet pendant cinq minutes : repérer les données, les documents, les barèmes et les questions indépendantes.
\item Répartir le temps selon le barème, environ 5 minutes par point : 40 minutes pour un exercice sur 8 points, 35 minutes pour
7 points, 25 minutes pour 5 points ; garder cinq à dix minutes pour relire.
\item Commencer par l'exercice que l'on maîtrise le mieux ; numéroter clairement chaque réponse.
\item Ne pas rester bloqué : une question « Montrer que » donne le résultat à utiliser pour la suite.
\end{questions}
\editionEtape{Comprendre les verbes de consigne}
\begin{tabularx}{\linewidth}{@{}p{38mm}Y@{}}\toprule
Calculer & obtenir une valeur numérique, avec son unité et un nombre de chiffres significatifs adapté. \\
Déterminer & obtenir un résultat par la méthode de son choix (calcul, lecture graphique), en l'explicitant. \\
Montrer que, établir & mener un raisonnement complet ; le résultat est connu, seul le raisonnement est évalué. \\
Exprimer & donner une expression littérale, en fonction des grandeurs indiquées. \\
Justifier & s'appuyer sur une loi, une propriété, une donnée ou un document. \\
Représenter, schématiser & tracer avec soin, légender, indiquer l'échelle d'un vecteur. \\
Interpréter, commenter & confronter un résultat à un modèle, à une référence ou au contexte, puis conclure. \\\bottomrule
\end{tabularx}
\end{blocEdition}

\editionPage{methodes-devoir-2}{Rédiger une réponse}{Calculs, prise d'initiative et relecture}{Présenter ce qui est évalué · Relire sa copie avant de la rendre}
\begin{blocEdition}{methodes-devoir-2}
\fontsize{10}{13.5}\selectfont
\editionEtape{Rédiger un calcul en cinq temps}
\begin{questions}
\item Annoncer la loi ou la relation utilisée : « D'après la loi de Beer-Lambert\ldots ».
\item Établir l'expression littérale de la grandeur cherchée, avant toute application numérique.
\item Remplacer par les valeurs, converties dans des unités cohérentes (le plus souvent celles du Système international).
\item Écrire le résultat avec son unité et un nombre adapté de chiffres significatifs, puis le mettre en évidence.
\item Conclure par une phrase qui répond à la question posée.
\end{questions}
\editionEtape{Réussir une question avec prise d'initiative}
Faire un schéma de la situation ; lister les données utiles des documents ; nommer la grandeur cherchée ; choisir les lois
qui la relient aux données ; mener le calcul ; vérifier l'ordre de grandeur et conclure. Une démarche cohérente, même
incomplète, rapporte des points : il faut donc l'écrire.
\editionEtape{Grille de relecture}
\begin{tabularx}{\linewidth}{@{}p{42mm}Y@{}}\toprule
\textbf{Attendu} & \textbf{Ce que le correcteur observe} \\\midrule
Rédaction & Phrases complètes et correctement construites ; accords ; ponctuation ; abréviations évitées. \\[1mm]
Raisonnement & Enchaînement logique explicite (« d'après », « or », « donc ») ; hypothèses énoncées ; conclusion qui répond à
la question. \\[1mm]
Vocabulaire & Termes précis : espèce et entité, transformation et réaction, oxydant et réducteur, masse et quantité de matière,
force et énergie, vitesse et célérité. \\[1mm]
Présentation & Questions numérotées ; résultats mis en évidence ; schémas légendés ; graphiques avec axes, grandeurs et
unités ; vecteurs avec leur échelle. \\\bottomrule
\end{tabularx}
\editionEtape{Les erreurs les plus coûteuses}
Un résultat sans unité ou avec trop de chiffres ; un oubli de conversion (mL en L, cm en m, g en kg, km/h en m/s, eV en J) ;
une confusion entre concentration en masse et concentration en quantité de matière ; un tableau d'avancement sans
réactif limitant ; un signe oublié en optique ; un résultat aberrant non commenté.
\editionRetenir{Une copie claire se lit vite : une idée par phrase, une loi citée avant chaque calcul, un résultat encadré avec
son unité, une conclusion qui reprend les mots de la question.}
\end{blocEdition}

\editionPage{capacites-experimentales}{Capacités expérimentales}{Ce qu'il faut savoir faire au laboratoire}{Programme de première · Renvois aux chapitres du livre}
\begin{blocEdition}{capacites-experimentales}
\fontsize{10}{13.5}\selectfont
Deux capacités sont communes à tous les thèmes : respecter les règles de sécurité liées au travail en laboratoire ; mettre
en œuvre un logiciel de simulation et de traitement des données (tableur, Python, interface d'acquisition).
\editionEtape{Constitution et transformations de la matière}
\begin{tabularx}{\linewidth}{@{}Yr@{}}\toprule
Préparer une solution par dissolution ou par dilution en choisissant le matériel adapté & ch. 1 \\
Réaliser le spectre d'absorption UV-visible d'une espèce ; mesurer des absorbances en s'aidant d'une notice & ch. 2 \\
Mettre en œuvre un test de reconnaissance pour identifier une espèce chimique & ch. 3 \\
Mettre en œuvre un titrage direct avec repérage colorimétrique de l'équivalence & ch. 5 \\
Visualiser la géométrie d'entités avec des modèles moléculaires ou un logiciel & ch. 6 \\
Proposer et réaliser une extraction liquide-liquide à partir de données de solubilité et de miscibilité & ch. 7 \\
Chauffage à reflux, distillation fractionnée, filtration, lavage, chromatographie sur couche mince & ch. 9 \\
Estimer une température de changement d'état ; éliminer les déchets de façon adaptée & ch. 9 \\\bottomrule
\end{tabularx}
\editionEtape{Mouvement, énergie, ondes et signaux}
\begin{tabularx}{\linewidth}{@{}Yr@{}}\toprule
Illustrer l'interaction électrostatique ; repérer la direction d'un champ électrique & ch. 11 \\
Mesurer une pression dans un gaz et dans un liquide & ch. 12 \\
Collecter des données sur un mouvement (vidéo, chronophotographie) & ch. 13 \\
Utiliser un multimètre ; réaliser un montage d'après un schéma normalisé ; microcontrôleur & ch. 14 \\
Estimer une énergie transférée électriquement ou mécaniquement & ch. 14 à 16 \\
Mesurer la période, la longueur d'onde, la célérité d'une onde périodique & ch. 17 \\
Estimer une distance focale ; former l'image d'un objet par une lentille & ch. 18 \\
Illustrer les synthèses additive et soustractive et la couleur apparente des objets & ch. 19 \\
Obtenir un spectre d'émission & ch. 20 \\
Exploiter une série de mesures ; évaluer une incertitude-type & ch. 21 \\\bottomrule
\end{tabularx}
\editionRetenir{Au laboratoire, une mesure n'est jamais seule : elle s'accompagne de son unité, de son incertitude-type et d'une
phrase qui la compare à la valeur attendue.}
\end{blocEdition}

\editionPage{fiches-experimentales}{Fiche 1 · La verrerie}{Fiches expérimentales : verrerie, montages, mesures}{Choisir le bon instrument · Préparer une solution}
\begin{blocEdition}{fiches-experimentales}
\fontsize{10}{13.5}\selectfont
\label{figedition-fiches-1}\schema{annexe-verrerie}
\editionEtape{Précision de la verrerie}
{\small\begin{tabularx}{\linewidth}{@{}lYl@{}}\toprule
\textbf{Verrerie} & \textbf{Usage} & \textbf{Tolérance usuelle (classe A)} \\\midrule
Fiole jaugée de \qty{100,0}{mL} & préparer un volume précis de solution & $\pm\qty{0,10}{mL}$ \\
Pipette jaugée de \qty{10,00}{mL} & prélever un volume précis & $\pm\qty{0,02}{mL}$ \\
Burette graduée de \qty{25}{mL} & verser un volume variable (titrage) & $\pm\qty{0,03}{mL}$, lecture à \qty{0,05}{mL} \\
Éprouvette graduée de \qty{100}{mL} & mesurer un volume approximatif & $\pm\qty{0,5}{mL}$ \\
Bécher, erlenmeyer & contenir, mélanger, chauffer & graduations indicatives \\\bottomrule
\end{tabularx}}\par\smallskip
L'incertitude-type associée à une tolérance $\pm a$ vaut $\dfrac{a}{\sqrt3}$ : \qty{0,058}{mL} pour la fiole de \qty{100,0}{mL}.
\editionEtape{Préparer une solution par dilution}
\begin{questions}
\item Calculer le volume de solution mère à prélever : $V_\text{mère}=\dfrac{c_\text{fille}\,V_\text{fille}}{c_\text{mère}}$ ; le facteur de dilution vaut
$F=\dfrac{c_\text{mère}}{c_\text{fille}}=\dfrac{V_\text{fille}}{V_\text{mère}}$.
\item Verser un peu de solution mère dans un bécher ; la prélever avec une pipette jaugée rincée avec cette solution.
\item La verser dans une fiole jaugée, compléter avec de l'eau distillée aux trois quarts, agiter.
\item Ajuster au trait de jauge (bas du ménisque, à hauteur des yeux), boucher et homogénéiser.
\end{questions}
\editionEtape{Préparer une solution par dissolution}
Peser la masse $m=c\,V\,M$ de soluté dans une coupelle, sur une balance tarée ; l'introduire dans la fiole jaugée avec un
entonnoir, rincer à l'eau distillée, dissoudre, compléter au trait de jauge et homogénéiser.
\end{blocEdition}

\editionPage{fiches-experimentales-2}{Fiche 2 · Doser une espèce en solution}{Fiches expérimentales : verrerie, montages, mesures}{Dosage par étalonnage · Titrage colorimétrique}
\begin{blocEdition}{fiches-experimentales-2}
\fontsize{10}{13.5}\selectfont
\label{figedition-fiches-2}\schema{05-montage}
\editionEtape{Mener un titrage colorimétrique}
Prélever la solution titrée à la pipette jaugée et la verser dans un erlenmeyer posé sur une feuille blanche ; remplir la burette
de solution titrante, chasser la bulle d'air, ajuster le zéro. Faire d'abord un titrage rapide, millilitre par millilitre, pour
encadrer l'équivalence, puis un titrage précis, goutte à goutte au voisinage de l'équivalence. L'équivalence est repérée par un
changement de couleur persistant (indicateur coloré, ou réactif lui-même coloré). À l'équivalence, pour la réaction
$a\,\mathrm{A}+b\,\mathrm{B}\to\ldots$ : $\dfrac{n_\mathrm{A}}{a}=\dfrac{n_{\mathrm{B},\mathrm{E}}}{b}$.
\editionEtape{Doser par étalonnage spectrophotométrique}
\begin{questions}
\item Faire le blanc avec le solvant ; régler le spectrophotomètre à la longueur d'onde du maximum d'absorption.
\item Mesurer l'absorbance de solutions étalons de concentrations connues ; tracer $A=f(c)$ : droite passant par l'origine (loi de
Beer-Lambert, $A=k\,c$, valable pour des solutions diluées).
\item Mesurer l'absorbance de la solution inconnue, diluée si nécessaire pour rester dans le domaine d'étalonnage ; en déduire
sa concentration par lecture graphique ou par le calcul, sans oublier le facteur de dilution.
\end{questions}
\end{blocEdition}

\editionPage{fiches-experimentales-3}{Fiche 3 · Synthèse et séparation}{Fiches expérimentales : verrerie, montages, mesures}{Chauffer, extraire, purifier, identifier}
\begin{blocEdition}{fiches-experimentales-3}
\fontsize{10}{13.5}\selectfont
\label{figedition-fiches-3}\schema{09-reflux}
\editionEtape{Les étapes d'une synthèse}
\begin{tabularx}{\linewidth}{@{}p{30mm}Y@{}}\toprule
Transformation & Chauffage à reflux : la température accélère la transformation, le réfrigérant condense les vapeurs, sans
perte de matière. Pierre ponce pour régulariser l'ébullition ; catalyseur éventuel. \\[1mm]
Séparation & Extraction liquide-liquide en ampoule à décanter : solvant non miscible à l'eau, dans lequel l'espèce est très
soluble ; la phase la moins dense est au-dessus. Filtration simple ou sous pression réduite pour un solide. \\[1mm]
Purification & Lavage, séchage par un desséchant anhydre ; distillation fractionnée d'un liquide. \\[1mm]
Identification & Température de fusion au banc Kofler ; chromatographie sur couche mince ; spectroscopie infrarouge. \\\bottomrule
\end{tabularx}
\editionEtape{Réaliser une chromatographie sur couche mince}
Déposer, sur la ligne de dépôt tracée au crayon, le produit obtenu et les références ; placer la plaque dans la cuve saturée
en vapeurs d'éluant, la ligne de dépôt au-dessus du niveau de l'éluant ; sortir la plaque avant que le front n'atteigne le
haut, marquer le front, sécher, révéler (lampe UV, diiode). Deux espèces identiques migrent à la même hauteur ; un produit
pur ne donne qu'une tache.
\editionEtape{Rendement et sécurité}
Rendement : $\eta=\dfrac{n_\text{obtenue}}{n_\text{max}}$, où $n_\text{max}$ se calcule à partir du réactif limitant. Porter blouse, lunettes et gants ;
manipuler les espèces volatiles ou corrosives sous la hotte ; consulter les pictogrammes de danger ; éliminer les déchets dans
les récipients prévus.
\end{blocEdition}

\editionPage{fiches-experimentales-4}{Fiche 4 · Mesures en physique}{Fiches expérimentales : verrerie, montages, mesures}{Acquérir, traiter, évaluer les incertitudes}
\begin{blocEdition}{fiches-experimentales-4}
\fontsize{10}{13.5}\selectfont
\editionEtape{Acquérir un signal}
\begin{tabularx}{\linewidth}{@{}p{38mm}Y@{}}\toprule
Multimètre & Brancher un voltmètre en dérivation, un ampèremètre en série ; commencer par le plus grand calibre, puis
choisir le plus petit calibre supérieur à la valeur mesurée. \\[1mm]
Oscilloscope, interface & Régler la base de temps pour voir quelques périodes ; mesurer une période sur plusieurs
périodes ; un retard entre deux signaux se mesure entre deux points homologues. \\[1mm]
Microcontrôleur & Commander un signal, lire un capteur (pression, température, distance), horodater les mesures. \\[1mm]
Pointage vidéo & Placer une échelle de longueur dans le plan du mouvement ; repérer l'origine des dates ; exporter les
coordonnées vers un tableur ou Python. \\\bottomrule
\end{tabularx}
\editionEtape{Optique}
Distance focale d'une lentille convergente : former sur un écran l'image nette d'un objet éloigné (elle se forme dans le plan
focal image), ou utiliser la relation de conjugaison $\dfrac{1}{\alg{OA'}}-\dfrac{1}{\alg{OA}}=\dfrac{1}{f'}$ pour plusieurs positions de l'objet.
Spectre d'émission : observer une lampe spectrale au spectroscope et repérer la longueur d'onde des raies.
\editionEtape{Évaluer une incertitude}
\begin{tabularx}{\linewidth}{@{}p{38mm}Y@{}}\toprule
Série de $n$ mesures (type A) & moyenne $\overline{x}$ et $u(\overline{x})=\dfrac{s}{\sqrt n}$, $s$ étant l'écart-type expérimental. \\[1mm]
Mesure unique (type B) & lecture sur une graduation : $u=\dfrac{\text{1 graduation}}{\sqrt{12}}$ ; précision $\pm a$ du fabricant : $u=\dfrac{a}{\sqrt3}$. \\[1mm]
Écriture du résultat & $u$ arrondie à un ou deux chiffres significatifs ; valeur arrondie au même rang ; unité. \\[1mm]
Comparaison & qualitative : un écart à la valeur de référence du même ordre que $u$ indique un résultat compatible. \\\bottomrule
\end{tabularx}
\editionRetenir{Mesurer plusieurs périodes, plusieurs longueurs d'onde ou plusieurs gouttes, puis diviser : c'est le moyen le plus
simple de réduire l'incertitude d'une mesure.}
\end{blocEdition}
